\documentclass[UTF8]{ctexart}
\usepackage[a4paper,margin=2.6cm]{geometry}
\usepackage{amsmath, amssymb, amsthm}
\usepackage{graphicx}
\usepackage[colorlinks=true,linkcolor=blue]{hyperref}

\title{复变函数学习笔记：Theta 函数的函数方程}
\author{}
\date{}

\newtheorem{theorem}{定理}[section]
\newtheorem{lemma}[theorem]{引理}
\newtheorem{proposition}[theorem]{命题}
\newtheorem{definition}[theorem]{定义}
\newtheorem{corollary}[theorem]{推论}
\newtheorem{remark}{注记}[section]
\numberwithin{equation}{section}

\newcommand{\RR}{\mathbb{R}}
\newcommand{\CC}{\mathbb{C}}
\newcommand{\ZZ}{\mathbb{Z}}
\begin{document}

\maketitle

本节从 Schwartz 空间（速降空间）中傅里叶变换的基本性质出发，严格证明泊松求和公式，并由此推出 Theta 函数所满足的函数方程 \(\vartheta(t) = t^{-1/2}\vartheta(1/t)\)。这一方程是第 6 章 Zeta 函数解析延拓的关键工具。

% ========== 1 ==========
\section{速降空间与傅里叶变换}

\begin{definition}[Schwartz 空间]
Schwartz 空间 \(\mathcal{S}(\RR)\) 是所有光滑函数 \(f:\RR\to\CC\) 的集合，满足对任意非负整数 \(\alpha,\beta\)，
\begin{equation}\label{eq:Schwartz}
\sup_{x\in\RR} |x|^\alpha |f^{(\beta)}(x)| < \infty .
\end{equation}
即函数及其各阶导数在无穷远处比任何多项式的倒数衰减得都快。
\end{definition}

\(\mathcal{S}(\RR)\) 是线性空间，且在逐点乘法、微分、傅里叶变换下封闭。傅里叶变换定义为
\begin{equation}
\hat{f}(\xi) = \int_{-\infty}^{\infty} f(x) e^{-2\pi i x\xi}\,dx ,\qquad \xi\in\RR .
\end{equation}
若 \(f\in\mathcal{S}(\RR)\)，则 \(\hat{f}\in\mathcal{S}(\RR)\)，且傅里叶反演公式成立：
\begin{equation}
f(x) = \int_{-\infty}^{\infty} \hat{f}(\xi) e^{2\pi i x\xi}\,d\xi ,\qquad x\in\RR .
\end{equation}
这一性质是证明泊松求和公式的核心。

% ========== 2 ==========
\section{泊松求和公式}

\begin{theorem}[泊松求和公式]\label{thm:Poisson}
设 \(f\in\mathcal{S}(\RR)\)，则
\begin{equation}\label{eq:Poisson}
\sum_{n=-\infty}^{\infty} f(n) = \sum_{n=-\infty}^{\infty} \hat{f}(n).
\end{equation}
\end{theorem}

\begin{proof}
由于 \(f\) 速降，级数 \(\sum_{n\in\ZZ} f(x+n)\) 在 \(\RR\) 上绝对且局部一致收敛，因此定义了周期为 \(1\) 的光滑函数
\[
F(x) = \sum_{n=-\infty}^{\infty} f(x+n).
\]
将周期函数 \(F\) 展开为傅里叶级数：
\[
F(x) = \sum_{k=-\infty}^{\infty} c_k e^{2\pi i k x},
\]
其中傅里叶系数
\[
c_k = \int_0^1 F(x) e^{-2\pi i k x}\,dx .
\]
将 \(F\) 的级数表示代入并交换积分与求和（由一致收敛保证）：
\begin{align*}
c_k &= \int_0^1 \sum_{n=-\infty}^{\infty} f(x+n) e^{-2\pi i k x}\,dx \\
&= \sum_{n=-\infty}^{\infty} \int_0^1 f(x+n) e^{-2\pi i k x}\,dx .
\end{align*}
对每个积分作变量代换 \(y = x+n\)，则
\[
\int_0^1 f(x+n) e^{-2\pi i k x}\,dx = \int_n^{n+1} f(y) e^{-2\pi i k (y-n)}\,dy = \int_n^{n+1} f(y) e^{-2\pi i k y}\,dy ,
\]
因为 \(e^{2\pi i k n}=1\)。将各区间上的积分求和，得到整个实轴上的积分：
\[
c_k = \int_{-\infty}^{\infty} f(y) e^{-2\pi i k y}\,dy = \hat{f}(k).
\]
因此
\[
F(x) = \sum_{k=-\infty}^{\infty} \hat{f}(k) e^{2\pi i k x}.
\]
特别地，取 \(x=0\)，即得
\[
\sum_{n=-\infty}^{\infty} f(n) = F(0) = \sum_{k=-\infty}^{\infty} \hat{f}(k),
\]
这就是泊松求和公式。
\end{proof}

% ========== 3 ==========
\section{Theta 函数的函数方程}

现在将泊松求和公式应用于特殊的高斯函数 \(f(x) = e^{-\pi t x^2}\)，其中 \(t>0\) 为固定参数。

\subsection*{高斯函数的傅里叶变换}

熟知对任意 \(a>0\)，
\[
\int_{-\infty}^{\infty} e^{-\pi a x^2} e^{-2\pi i x\xi}\,dx = \frac{1}{\sqrt{a}} e^{-\pi \xi^2 / a}.
\]
取 \(a = t\)，则 \(f(x) = e^{-\pi t x^2}\) 的傅里叶变换为
\begin{equation}\label{eq:GaussFT}
\hat{f}(\xi) = \frac{1}{\sqrt{t}} e^{-\pi \xi^2 / t}.
\end{equation}

\subsection*{应用泊松求和}

显然 \(f\in\mathcal{S}(\RR)\)。由定理 \ref{thm:Poisson}，
\[
\sum_{n=-\infty}^{\infty} f(n) = \sum_{n=-\infty}^{\infty} \hat{f}(n).
\]
左边是 \(\sum_{n=-\infty}^{\infty} e^{-\pi t n^2}\)，右边代入 \eqref{eq:GaussFT} 得 \(\sum_{n=-\infty}^{\infty} \frac{1}{\sqrt{t}} e^{-\pi n^2 / t}\)。因此
\[
\sum_{n=-\infty}^{\infty} e^{-\pi t n^2} = \frac{1}{\sqrt{t}} \sum_{n=-\infty}^{\infty} e^{-\pi n^2 / t}.
\]

\subsection*{Theta 函数表示}

定义 Theta 函数（本书第 6 章所用形式）为
\[
\vartheta(t) = \sum_{n=-\infty}^{\infty} e^{-\pi n^2 t}, \qquad t>0.
\]
则上面的恒等式正是
\begin{equation}\label{eq:ThetaFuncEq}
\boxed{\vartheta(t) = t^{-1/2} \vartheta(1/t)} .
\end{equation}
这就是 Theta 函数所满足的函数方程，它是 Riemann Zeta 函数解析延拓的出发点。

\end{document}